\chapter{THIẾT KẾ KIẾN TRÚC, DỮ LIỆU VÀ TÍCH HỢP AI}

\section{Kiến trúc tổng thể 4 tầng (C4 Container)}

Hệ thống \textbf{GarageTools AI} được tổ chức theo kiến trúc 4 tầng tách biệt rõ trách nhiệm: (1) Tầng Giao diện người dùng (\texttt{frontend/src/} chạy React 19 + TanStack Start song song cùng giao diện Jinja2 tại \texttt{app/templates/}); (2) Tầng Định tuyến \& Bảo mật (\texttt{app/main.py}, \texttt{app/core/}, \texttt{app/routes/}); (3) Tầng Dịch vụ Nghiệp vụ \& AI Gateway (\texttt{app/services/}, \texttt{app/ai\_gateway.py}, \texttt{config/prompts/}); và (4) Tầng Cơ sở dữ liệu quan hệ SQLite WAL (\texttt{garage.db}). Hình~\ref{fig:g-container} minh họa sơ đồ kiến trúc thành phần chứa của hệ thống.

\begin{figure}[H]
\centering
\vspace{2pt}
\resizebox{\textwidth}{!}{%
\begin{tikzpicture}[
    font=\sffamily\footnotesize,
    >=Latex,
    cardblue/.style={rectangle, draw=diagblueline, line width=0.9pt, fill=diagbluefill, rounded corners=5pt, text width=4.3cm, align=center, inner sep=5.5pt},
    cardgreen/.style={rectangle, draw=diaggreenline, line width=0.9pt, fill=diaggreenfill, rounded corners=5pt, text width=4.3cm, align=center, inner sep=5.5pt},
    cardamber/.style={rectangle, draw=diagamberline, line width=0.9pt, fill=diagamberfill, rounded corners=5pt, text width=4.3cm, align=center, inner sep=5.5pt},
    cardpurple/.style={rectangle, draw=diagpurpleline, line width=0.9pt, fill=diagpurplefill, rounded corners=5pt, text width=4.3cm, align=center, inner sep=5.5pt},
    layerbar/.style={rectangle, draw=diagslateline, line width=0.8pt, fill=diagslatefill, rounded corners=4pt, text width=14.5cm, align=center, font=\sffamily\scriptsize\bfseries, inner sep=4pt},
    arrow/.style={-Latex, line width=0.9pt, diagblueline}
]

% Tầng 1: Giao diện
\node[cardblue, text width=14.5cm] (fe) {
    \textbf{TẦNG 1: GIAO DIỆN NGƯỜI DÙNG (\texttt{frontend/src/} \& \texttt{app/templates/})}\\
    React 19 + TanStack Start + TypeScript + Tailwind CSS v4 (16 màn hình nghiệp vụ kết nối \texttt{/api/*} kèm Jinja2 SSR)
};

% Tầng 2: Định tuyến & Bảo mật (dùng below=0.45cm of ... để không bao giờ đè nhau)
\node[layerbar, below=0.45cm of fe] (bar2) {TẦNG 2: ĐỊNH TUYẾN, XÁC THỰC PHIÊN \& MA TRẬN PHÂN QUYỀN (\texttt{app/main.py}, \texttt{app/core/}, \texttt{app/routes/})};

\node[cardpurple, below=0.35cm of bar2] (rbac) {
    \textbf{Ma trận Phân quyền RBAC}\\
    4 vai trò (\texttt{security\_matrix.py})\\
    Cô lập phiếu KTV (\texttt{duoc\_gan\_cong})
};
\node[cardpurple, left=0.5cm of rbac] (sec) {
    \textbf{Bảo mật \& Giới hạn Tần suất}\\
    Argon2id + Cookie \texttt{garage\_session}\\
    CSRF 2 lớp + \texttt{auth\_version}
};
\node[cardpurple, right=0.5cm of rbac] (api) {
    \textbf{Hệ thống Route \& JSON API}\\
    14 mô-đun JSON (\texttt{routes/api/})\\
    + 7 mô-đun Jinja2 (\texttt{routes/})
};

% Tầng 3: Dịch vụ nghiệp vụ & AI
\node[layerbar, below=0.5cm of rbac] (bar3) {TẦNG 3: DỊCH VỤ NGHIỆP VỤ LÕI \& CỔNG TRỢ LÝ AI (\texttt{app/services/}, \texttt{app/ai\_gateway.py}, \texttt{config/prompts/})};

\node[cardamber, below=0.35cm of bar3] (aigw) {
    \textbf{Cổng AI (\texttt{ai\_gateway.py})}\\
    AI01 \texttt{history}, AI02 \texttt{explain}, AI03 \texttt{quote}\\
    Lọc khớp giá CSDL + ép \texttt{backend\_total}
};
\node[cardgreen, left=0.5cm of aigw] (svc) {
    \textbf{11 Mô-đun Nghiệp vụ Lõi}\\
    \texttt{orders}, \texttt{quotations}, \texttt{stock},\\
    \texttt{billing}, \texttt{warranties}, \texttt{parties}...
};
\node[cardamber, right=0.5cm of aigw] (prompt) {
    \textbf{Tệp Prompt Ngoại hóa}\\
    \texttt{system.txt}, \texttt{history.txt},\\
    \texttt{explain.txt}, \texttt{quote.txt}
};

% Tầng 4: CSDL & Cloud AI
\node[cardblue, text width=6.9cm, below=0.6cm of $(svc.south)!0.5!(aigw.south)$, xshift=-1.1cm] (db) {
    \textbf{TẦNG 4A: CSDL QUAN HỆ SQLITE 3 (\texttt{garage.db})}\\
    Chế độ WAL + \texttt{PRAGMA foreign\_keys=ON} + \texttt{BEGIN IMMEDIATE}\\
    21 bảng chuẩn 3NF, 37 chỉ mục, 171 dòng dữ liệu mẫu
};
\node[cardgreen, text width=6.9cm, right=0.5cm of db] (cloud) {
    \textbf{TẦNG 4B: NHÀ CUNG CẤP MÔ HÌNH AI (7 PRESETS)}\\
    \texttt{groq} (\texttt{openai/gpt-oss-120b}), \texttt{gemini} (\texttt{gemini-2.0-flash}),\\
    \texttt{openrouter}, \texttt{openai}, \texttt{9router}, \texttt{custom} \& \texttt{mock}
};

% Kết nối giữa các tầng
\draw[arrow] (fe.south) -- (bar2.north);
\draw[arrow] (sec.east) -- (rbac.west);
\draw[arrow] (rbac.east) -- (api.west);
\draw[arrow] (rbac.south) -- (bar3.north);
\draw[arrow, <->] (aigw.east) -- (prompt.west);
\draw[arrow, <->] (svc.south) -- (svc.south |- db.north);
\draw[arrow, <->] ([xshift=-0.8cm]aigw.south) -- ([xshift=2.2cm]db.north);
\draw[arrow, <->, diaggreenline] ([xshift=0.8cm]aigw.south) -- ([xshift=-2.2cm]cloud.north);

\end{tikzpicture}%
}

\vspace{2pt}
\caption{Sơ đồ kiến trúc thành phần chứa (C4 Container) của hệ thống GarageTools AI}
\label{fig:g-container}
\end{figure}

Mọi yêu cầu từ trình duyệt đều đi qua lớp kiểm tra phiên đăng nhập (\texttt{garage\_session}), đối chiếu \texttt{auth\_version}, kiểm tra CSRF token (\texttt{garage\_csrf} / \texttt{X-CSRF-Token}) và kiểm tra quyền trong \texttt{security\_matrix.py} trước khi thực thi logic tại \texttt{app/services/}.

\section{Thiết kế cơ sở dữ liệu quan hệ (ERD 21 bảng chuẩn 3NF)}

Cơ sở dữ liệu của hệ thống gồm \textbf{21 bảng} (17 bảng nghiệp vụ lõi của garage và 4 bảng mở rộng cửa hàng dụng cụ \texttt{shop\_*}) cùng \textbf{37 chỉ mục (indexes)}, được định nghĩa tại \texttt{database/schema.sql} và tự động nâng cấp qua 5 bản di trú (\texttt{001} $\rightarrow$ \texttt{005}) trong \texttt{database/migrations/}. Sơ đồ ERD được chia thành hai nhóm: Phân hệ Tiếp nhận, Phiếu sửa chữa \& Sổ cái Kho (Hình~\ref{fig:g-erd-core}) và Phân hệ Người dùng, Báo giá, Hóa đơn, Bảo hành \& Kết quả AI (Hình~\ref{fig:g-erd-ops}).

\begin{figure}[H]
\centering
\vspace{2pt}
\resizebox{\textwidth}{!}{%
\begin{tikzpicture}[
    font=\sffamily\scriptsize,
    >=Latex,
    entity/.style={rectangle, draw=diagblueline, line width=0.9pt, fill=white, rounded corners=4pt, inner sep=0pt},
    rel/.style={line width=0.9pt, diagblueline, -Latex}
]

\node[entity] (cust) at (0,0) {
\begin{tabular}{p{4.3cm}}
\rowcolor{diagbluefill}\textbf{1. customers}\\
\hline
\textbf{PK} id \textit{INTEGER}\\
name, phone \textit{UNIQUE}\\
email, address, created\_at\\
\end{tabular}
};

\node[entity] (veh) at (6.2,0) {
\begin{tabular}{p{4.3cm}}
\rowcolor{diagbluefill}\textbf{2. vehicles}\\
\hline
\textbf{PK} id, \textbf{FK} customer\_id\\
plate \textit{UNIQUE}, brand, model\\
year \textit{1900..2100}, vin \textit{UNIQUE}\\
odometer \textit{>=0}, is\_active\\
\end{tabular}
};

\node[entity] (appt) at (12.4,0) {
\begin{tabular}{p{4.3cm}}
\rowcolor{diagbluefill}\textbf{3. appointments}\\
\hline
\textbf{PK} id, \textbf{FK} customer\_id, vehicle\_id\\
scheduled\_at, notes\\
status (\textit{4 trạng thái}), created\_at\\
\end{tabular}
};

\node[entity, draw=diaggreenline] (roi) at (0,-4.2) {
\begin{tabular}{p{4.3cm}}
\rowcolor{diaggreenfill}\textbf{4. repair\_order\_items \& logs}\\
\hline
\textbf{PK} id, \textbf{FK} repair\_order\_id\\
\textit{CHECK XOR}(service\_id, part\_id)\\
qty \textit{>0}, unit\_price \textit{>=0 (chốt cứng)}\\
\textbf{1-N} repair\_logs (technician\_id, note)\\
\end{tabular}
};

\node[entity, draw=diagamberline] (ro) at (6.2,-4.2) {
\begin{tabular}{p{4.3cm}}
\rowcolor{diagamberfill}\textbf{5. repair\_orders \& intake\_records}\\
\hline
\textbf{PK} id, \textbf{FK} vehicle\_id, technician\_id\\
status (\textit{6 trạng thái}), technical\_notes\\
acceptance\_note, accepted\_by, acceptance\_at\\
\textbf{1-1} intake\_records (ODO, nhiên liệu, ngoại thất)\\
\end{tabular}
};

\node[entity, draw=diagpurpleline] (stock) at (12.4,-4.2) {
\begin{tabular}{p{4.3cm}}
\rowcolor{diagpurplefill}\textbf{6. services, parts \& stock\_txns}\\
\hline
\textbf{services}: id, name \textit{UNIQUE}, price\\
\textbf{parts}: id, name \textit{UNIQUE}, sku, stock\\
\textbf{stock\_transactions} (Sổ cái kho):\\
reason \textit{RECEIPT|USAGE|RETURN|ADJ}, qty\\
\end{tabular}
};

\draw[rel] (cust.east) -- node[above, font=\scriptsize]{1 : N} (veh.west);
\draw[rel] (veh.east) -- node[above, font=\scriptsize]{1 : N} (appt.west);
\draw[rel] (veh.south) -- node[right, font=\scriptsize]{1 : N} (ro.north);
\draw[rel] (ro.west) -- node[above, font=\scriptsize]{1 : N} (roi.east);
\draw[rel] (stock.west) -- node[above, font=\scriptsize]{1 : N} (ro.east);

\end{tikzpicture}%
}

\vspace{2pt}
\caption{Sơ đồ thực thể liên kết (ERD) phân hệ Khách hàng, Phương tiện, Phiếu sửa chữa và Sổ cái Kho}
\label{fig:g-erd-core}
\end{figure}

\begin{figure}[H]
\centering
\vspace{2pt}
\resizebox{0.98\textwidth}{!}{%
\begin{tikzpicture}[
  font=\sffamily,
  ent/.style={
    rectangle split,
    rectangle split parts=2,
    rectangle split part fill={black!12,white},
    draw=black,
    line width=0.65pt,
    text width=3.9cm,
    align=left,
    font=\sffamily\scriptsize,
    inner sep=4pt
  },
  rel/.style={-{Latex[length=1.8mm,width=1.4mm]}, line width=0.6pt, draw=black},
  lbl/.style={font=\sffamily\scriptsize, fill=white, inner sep=1.2pt}
]
  \node[ent] (usr) at (0, 1.85) {
    \textbf{users (Tài khoản \& Vai trò)}
    \nodepart{second}
    \underline{id} (PK), username, email\\
    password\_hash, role (4 vai trò)\\
    is\_active, created\_at
  };

  \node[ent] (inv) at (5.6, 1.85) {
    \textbf{invoices \& payments}
    \nodepart{second}
    \underline{id} (PK), repair\_order\_id (FK)\\
    subtotal, discount, tax\_vat\\
    total\_payable, paid\_amount, status
  };

  \node[ent] (war) at (11.2, 1.85) {
    \textbf{warranty\_claims \& reminders}
    \nodepart{second}
    \underline{id} (PK), repair\_order\_id (FK)\\
    vehicle\_id (FK), due\_mileage\\
    due\_date, claim\_status
  };

  \node[ent] (aisug) at (0, -1.85) {
    \textbf{ai\_suggestions (Đề xuất HITL)}
    \nodepart{second}
    \underline{id} (PK), repair\_order\_id (FK)\\
    task\_code (AI01..AI04), payload\_json\\
    confidence, review\_status (HITL)\\
    reviewed\_by (FK), reviewed\_at
  };

  \node[ent] (ailog) at (5.6, -1.85) {
    \textbf{ai\_audit\_logs \& prompts}
    \nodepart{second}
    \underline{id} (PK), user\_id (FK), task\_code\\
    provider, model, prompt\_version\\
    latency\_ms, is\_fallback, status
  };

  \node[ent] (cred) at (11.2, -1.85) {
    \textbf{ai\_provider\_credentials}
    \nodepart{second}
    \underline{id} (PK), owner\_user\_id (FK)\\
    provider\_name, encrypted\_api\_key\\
    base\_url, is\_valid, updated\_at
  };

  \draw[rel] (usr.east) -- node[lbl, above]{1..n} (inv.west);
  \draw[rel] (inv.east) -- node[lbl, above]{1..n} (war.west);
  \draw[rel] (usr.south) -- node[lbl, right=2pt]{1..n} (aisug.north);
  \draw[rel] (aisug.east) -- node[lbl, above]{1..1} (ailog.west);
  \draw[rel] (ailog.east) -- node[lbl, above]{1..n} (cred.west);
\end{tikzpicture}%
}

\vspace{2pt}
\caption{Sơ đồ thực thể liên kết (ERD) phân hệ Người dùng, Báo giá, Hóa đơn, Bảo hành và Kết quả AI}
\label{fig:g-erd-ops}
\end{figure}

Bảng~\ref{tab:g-db-schema} đặc tả chính xác tên bảng, khóa chính, khóa ngoại và các ràng buộc toàn vẹn (\texttt{CHECK}, \texttt{UNIQUE}, \texttt{ON DELETE RESTRICT}) đúng như trong tệp \texttt{database/schema.sql}.

\begin{small}
\renewcommand{\arraystretch}{1.05}
\begin{longtable}{|L{4.2cm}|L{3.1cm}|L{6.4cm}|}
\caption{Đặc tả cấu trúc 21 bảng dữ liệu và ràng buộc toàn vẹn trong \texttt{database/schema.sql}}\label{tab:g-db-schema}\\
\hline
\thead{Tên bảng} & \thead{Khóa chính \& Khóa ngoại} & \thead{Vai trò nghiệp vụ và Ràng buộc toàn vẹn (Constraints)}\\
\hline
\endfirsthead

\hline
\thead{Tên bảng} & \thead{Khóa chính \& Khóa ngoại} & \thead{Vai trò nghiệp vụ và Ràng buộc toàn vẹn (Constraints)}\\
\hline
\endhead

\hline
\endfoot

\hline
\endlastfoot

1. \texttt{users} & PK: \texttt{id} & Tài khoản nhân viên; \texttt{username UNIQUE}; \texttt{role IN ('MANAGER', 'RECEPTIONIST', 'TECHNICIAN', 'CASHIER')}; \texttt{is\_active IN (0,1)}; \texttt{auth\_version >= 1}. \\
\hline
2. \texttt{customers} & PK: \texttt{id} & Hồ sơ khách hàng; lưu \texttt{name}, \texttt{phone UNIQUE}, \texttt{email}, \texttt{address}, \texttt{created\_at}. \\
\hline
3. \texttt{vehicles} & PK: \texttt{id}; FK: \texttt{customer\_id} & Hồ sơ phương tiện; \texttt{plate UNIQUE}, \texttt{vin UNIQUE}; \texttt{year BETWEEN 1900 AND 2100}; \texttt{odometer >= 0}; \texttt{is\_active IN (0,1)}. \\
\hline
4. \texttt{appointments} & PK: \texttt{id}; FK: \texttt{customer\_id}, \texttt{vehicle\_id} & Lịch hẹn dịch vụ; giữ cả \texttt{customer\_id} và \texttt{vehicle\_id} để đặt lịch nhanh; \texttt{status IN ('SCHEDULED', 'CONFIRMED', 'CANCELLED', 'DONE')}. \\
\hline
5. \texttt{services} & PK: \texttt{id} & Danh mục dịch vụ sửa chữa; \texttt{name UNIQUE}, \texttt{description} (cho phép \texttt{NULL}), \texttt{price >= 0} (kiểu \texttt{INTEGER}), \texttt{is\_active}. \\
\hline
6. \texttt{parts} & PK: \texttt{id} & Danh mục phụ tùng; \texttt{name UNIQUE}, \texttt{sku}, \texttt{stock >= 0}, \texttt{min\_stock >= 0}, \texttt{price >= 0}, \texttt{is\_active}. \\
\hline
7. \texttt{stock\_transactions} & PK: \texttt{id}; FK: \texttt{part\_id}, \texttt{repair\_order\_id}, \texttt{created\_by} & Sổ cái tồn kho bất biến; \texttt{reason IN ('RECEIPT', 'USAGE', 'RETURN', 'ADJUSTMENT')}; duy trì bất biến $\texttt{parts.stock} = \sum \texttt{qty}$. \\
\hline
8. \texttt{repair\_orders} & PK: \texttt{id}; FK: \texttt{vehicle\_id}, \texttt{technician\_id}, \texttt{accepted\_by} & Phiếu sửa chữa; 6 trạng thái (\texttt{RECEIVED}, \texttt{INSPECTING}, \texttt{ESTIMATING}, \texttt{REPAIRING}, \texttt{COMPLETED}, \texttt{CANCELLED}); \texttt{CHECK (acceptance\_at IS NULL OR accepted\_by IS NOT NULL)}. \\
\hline
9. \texttt{repair\_order\_items} & PK: \texttt{id}; FK: \texttt{repair\_order\_id}, \texttt{service\_id}, \texttt{part\_id} & Hạng mục của phiếu; \textbf{ràng buộc XOR}: \texttt{(service\_id IS NOT NULL AND part\_id IS NULL) OR (service\_id IS NULL AND part\_id IS NOT NULL)}; \texttt{qty > 0}, \texttt{unit\_price >= 0} (chốt cứng). \\
\hline
10. \texttt{repair\_logs} & PK: \texttt{id}; FK: \texttt{repair\_order\_id}, \texttt{technician\_id} & Nhật ký kỹ thuật ghi lại từng bước kiểm tra, sửa chữa và chuyển trạng thái trên phiếu. \\
\hline
11. \texttt{intake\_records} & PK: \texttt{id}; FK: \texttt{repair\_order\_id UNIQUE}, \texttt{created\_by} & Biên bản tiếp nhận xe (1--1 với phiếu sửa chữa); lưu \texttt{odometer\_at\_intake >= 0}, \texttt{fuel\_level}, \texttt{exterior\_condition}, \texttt{accessories}. \\
\hline
12. \texttt{quotations} & PK: \texttt{id}; FK: \texttt{repair\_order\_id}, \texttt{created\_by}, \texttt{approved\_by}, \texttt{applied\_by} & Báo giá chính thức; 5 trạng thái (\texttt{DRAFT}, \texttt{SENT}, \texttt{APPROVED}, \texttt{REJECTED}, \texttt{EXPIRED}); cột \texttt{applied\_at} đảm bảo chỉ áp dụng 1 lần (\textit{idempotent}). \\
\hline
13. \texttt{quotation\_items} & PK: \texttt{id}; FK: \texttt{quotation\_id}, \texttt{service\_id}, \texttt{part\_id} & Hạng mục trong báo giá; áp dụng cùng \textbf{ràng buộc XOR} giữa dịch vụ và phụ tùng; \texttt{qty > 0}, \texttt{unit\_price >= 0}. \\
\hline
14. \texttt{invoices} & PK: \texttt{id}; FK: \texttt{repair\_order\_id UNIQUE} & Hóa đơn của phiếu (1--1); \texttt{status IN ('UNPAID','PARTIAL','PAID')}; \textbf{ràng buộc chống thu vượt}: \texttt{CHECK (paid <= total)}. \\
\hline
15. \texttt{payments} & PK: \texttt{id}; FK: \texttt{invoice\_id} & Lịch sử các đợt thanh toán; \texttt{amount > 0}; \texttt{method IN ('CASH','CARD','TRANSFER','ZALO')}. \\
\hline
16. \texttt{ai\_results} & PK: \texttt{id}; FK: \texttt{repair\_order\_id}, \texttt{reviewed\_by} & Lưu kết quả AI và trạng thái duyệt; \texttt{kind IN ('history', 'explain', 'quote')}; \texttt{review\_status IN ('PENDING', 'APPROVED', 'REJECTED')}; \texttt{CHECK (review\_status = 'PENDING' OR (reviewed\_by IS NOT NULL AND reviewed\_at IS NOT NULL))}. \\
\hline
17. \texttt{warranties} & PK: \texttt{id}; FK: \texttt{repair\_order\_id}, \texttt{service\_id}, \texttt{part\_id}, \texttt{created\_by} & Bảo hành theo hạng mục của phiếu đã hoàn thành; \texttt{scope IN ('SERVICE', 'PART')} kèm \textbf{ràng buộc XOR}; \texttt{months > 0}, \texttt{end\_date > start\_date}. \\
\hline
18--21. \texttt{shop\_categories},\newline \texttt{shop\_products},\newline \texttt{shop\_orders},\newline \texttt{shop\_order\_items} & PK: \texttt{id}; FK: \texttt{category\_id}, \texttt{order\_id}, \texttt{product\_id} & Bốn bảng mở rộng cho mô-đun cửa hàng dụng cụ/phụ kiện (\texttt{shop\_orders.status IN ('PENDING', 'PAID', 'SHIPPED', 'CANCELLED')}). \\
\hline
\end{longtable}
\end{small}

\section{Thiết kế Pipeline AI 7 bước và Cơ chế Chống ảo giác giá}

Để đảm bảo AI hoạt động đúng vai trò trợ lý tư vấn, hệ thống thiết kế luồng xử lý 7 bước tại \texttt{app/ai\_gateway.py} và \texttt{app/services/ai.py} (Hình~\ref{fig:g-ai-pipeline}) cùng bảng kiểm soát dữ liệu cho từng chức năng (Bảng~\ref{tab:g-context-gate}).

\begin{figure}[H]
\centering
\vspace{2pt}
\resizebox{\textwidth}{!}{%
\begin{tikzpicture}[
    font=\sffamily\footnotesize,
    >=Latex,
    stepbox/.style={rectangle, draw=diagblueline, line width=0.9pt, fill=diagbluefill, rounded corners=5pt, text width=3.1cm, minimum height=1.65cm, align=center, inner sep=4.5pt},
    guardbox/.style={rectangle, draw=diagamberline, line width=0.9pt, fill=diagamberfill, rounded corners=5pt, text width=3.1cm, minimum height=1.65cm, align=center, inner sep=4.5pt},
    extbox/.style={rectangle, draw=diaggreenline, line width=0.9pt, fill=diaggreenfill, rounded corners=5pt, text width=3.1cm, minimum height=1.65cm, align=center, inner sep=4.5pt},
    arr/.style={-Latex, line width=0.9pt, diagblueline}
]

\node[stepbox] (s1) {\textbf{1. Cookie + CSRF}\\Kiểm tra quyền RBAC\\+ Giới hạn 10 lần/60s};
\node[stepbox, right=0.85cm of s1] (s2) {\textbf{2. Truy vấn CSDL}\\Đọc \texttt{vehicles}, \texttt{services},\\\texttt{repair\_order\_items}};
\node[guardbox, right=0.85cm of s2] (s3) {\textbf{3. Ghép Prompt}\\Nạp 4 tệp mẫu trong\\\texttt{config/prompts/}};
\node[extbox, right=0.85cm of s3] (s4) {\textbf{4. Gọi Mô hình AI}\\Groq / Gemini / OpenAI\\(Timeout 15 giây)};

\node[guardbox, below=0.65cm of s4] (s5) {\textbf{5. Lọc \& Dự phòng}\\Khớp đúng Tên + Đơn giá\\Ép \texttt{backend\_total}};
\node[extbox, below=0.65cm of s1] (s7) {\textbf{7. Human Review}\\Quản lý (\texttt{ai.review})\\Duyệt / Từ chối 1 lần};
\node[stepbox] (s6) at ($(s7)!0.5!(s5)$) {\textbf{6. Lưu \texttt{ai\_results}}\\Gán \texttt{PENDING}\\Lưu dữ liệu đầu vào/ra};

\draw[arr] (s1) -- (s2);
\draw[arr] (s2) -- (s3);
\draw[arr] (s3) -- (s4);
\draw[arr] (s4) -- (s5);
\draw[arr] (s5) -- (s6);
\draw[arr] (s6) -- (s7);

\end{tikzpicture}%
}

\vspace{2pt}
\caption{Quy trình xử lý 7 bước của AI Gateway, lọc khớp giá CSDL và duyệt Human Review}
\label{fig:g-ai-pipeline}
\end{figure}

\begin{table}[H]
\centering
\small
\caption{Dữ liệu đầu vào và cơ chế kiểm soát đầu ra của 3 chức năng AI}
\label{tab:g-context-gate}
\begin{tabularx}{\textwidth}{|L{2.9cm}|Y|Y|}
\hline
\thead{Chức năng AI} & \thead{Biến dữ liệu ghép vào Prompt (\texttt{config/prompts/*.txt})} & \thead{Ràng buộc đầu ra \& Xử lý ca biên tại \texttt{ai\_gateway.py}}\\
\hline
AI01: Tóm tắt lịch sử sửa chữa (\texttt{history.txt}) & Biến \texttt{\{\{repair\_history\}\}}: thông tin biển số, hãng, dòng, năm sản xuất và danh sách các phiếu sửa chữa + hạng mục trước đó của xe. & Trả về đúng 4 mục đánh số \texttt{1..4}; nếu xe chưa từng sửa chữa (\texttt{orders} rỗng, ví dụ \texttt{vehicle\_id=6}), ghi rõ “Chưa có lịch sử sửa chữa” và cấm suy đoán lỗi. \\
\hline
AI02: Giải thích dịch vụ (\texttt{explain.txt}) & Biến \texttt{\{\{service\}\}}: tên dịch vụ, đơn giá niêm yết (\texttt{services.price}) và mô tả kỹ thuật (\texttt{services.description}). & Trả về 3 mục dễ hiểu cho khách hàng; nếu dịch vụ chưa có mô tả (\texttt{description IS NULL}, ví dụ \texttt{service\_id=7}) thì tự điền câu mặc định ở mục 1, không báo lỗi 500. \\
\hline
AI03: Báo giá nháp (\texttt{quote.txt}) & Biến \texttt{\{\{items\}\}} và \texttt{\{\{backend\_total\}\}}: danh sách hạng mục từ \texttt{repair\_order\_items} (hoặc từ báo giá \texttt{quotations} đang mở) và tổng tiền do Backend tự cộng. & Chỉ giữ lại các dòng trong JSON khớp chính xác \texttt{(name, qty, unit\_price)} với CSDL; luôn gán \texttt{total = backend\_total}; tự động fallback sang \texttt{\_goi\_mock} nếu lỗi. \\
\hline
\end{tabularx}
\end{table}

\section{Ba vòng thử nghiệm và tối ưu Prompt (\texttt{KT3\_PROMPT\_EXPERIMENTS.md})}

Toàn bộ lời nhắc gửi lên mô hình được tách khỏi mã nguồn Python thành 4 tệp văn bản trong \texttt{config/prompts/} (\texttt{system.txt}, \texttt{history.txt}, \texttt{explain.txt}, \texttt{quote.txt}). Nhóm đã thực hiện 3 vòng thử nghiệm và cải tiến cho từng chức năng:

\begin{itemize}[noitemsep,topsep=2pt]
  \item \textbf{Vòng 1 -- Câu lệnh ngắn, không quy định cấu trúc đầu ra}:
  Ví dụ với AI01 chỉ viết \textit{“Tóm tắt lịch sử sửa chữa sau cho nhân viên lễ tân: \{\{repair\_history\}\}”}; với AI03 chỉ yêu cầu \textit{“Tạo báo giá nháp dạng JSON từ danh sách hạng mục”}.
  \textbf{Vấn đề gặp phải}: Mỗi lần gọi AI01 trả về một bố cục khác nhau, và khi gặp xe mới chưa có lịch sử thì mô hình tự suy diễn các lỗi thường gặp của dòng xe; với AI03, mô hình trả kèm khối \texttt{```json} và có lần cộng nhầm tổng tiền \texttt{total}.
  \item \textbf{Vòng 2 -- Quy định rõ bố cục mục và khóa JSON, truyền sẵn \texttt{backend\_total}}:
  Quy định cố định 4 mục cho AI01, 3 mục cho AI02 và truyền sẵn biến \texttt{\{\{backend\_total\}\}} vào \texttt{quote.txt}.
  \textbf{Kết quả \& Điểm còn thiếu}: Đầu ra đã đi vào khuôn khổ, nhưng khi kiểm thử với dịch vụ có \texttt{description IS NULL} (dịch vụ \texttt{\#7} cân chỉnh thước lái) ở AI02 thì mô hình giải thích chung chung và thiếu đơn giá niêm yết; với AI03 khi phiếu sửa chữa trống hạng mục thì chưa có quy định trả mảng rỗng rõ ràng.
  \item \textbf{Vòng 3 -- Bổ sung quy tắc xử lý dữ liệu trống và bộ lọc hậu kiểm ở Backend}:
  Hoàn thiện cả 4 tệp prompt trong \texttt{config/prompts/}: cấm suy đoán khi xe chưa có lịch sử, truyền kèm đơn giá niêm yết và quy định câu trả lời mặc định khi \texttt{description} trống cho AI02, đồng thời bổ sung bộ lọc đối chiếu từng dòng \texttt{(name, qty, unit\_price)} và cưỡng bức \texttt{total = backend\_total} trong hàm \texttt{bao\_gia\_nhap}.
\end{itemize}

\begin{lstlisting}[caption={Nội dung thực tế của tệp config/prompts/system.txt và config/prompts/quote.txt}]
# File: config/prompts/system.txt
Ban la tro ly AI noi bo cua garage sua chua o to. Nguyen tac bat buoc:
1. Chi su dung thong tin co trong du lieu duoc cung cap. Khong suy doan, khong bia
   them hang muc, gia tien, ngay thang hay nguyen nhan hong hoc.
2. Neu du lieu thieu hoac trong, noi ro la chua co thong tin; khong tu dien.
3. Khong dua ra chan doan ky thuat thay cho ky thuat vien.
4. Tra loi bang tieng Viet ro rang, ngan gon, dung trong tam.

# File: config/prompts/quote.txt
Lap bao gia nhap tu danh sach hang muc da duoc ghi nhan trong phieu sua chua.
Quy tac bat buoc:
- Chi dung cac hang muc co trong danh sach. Khong them hang muc moi, khong doi
  so luong, khong doi don gia.
- Giu nguyen tong tien do he thong tinh san: {{backend_total}} VND.
- Neu danh sach hang muc trong, tra ve "items": [] va "total": 0...
\end{lstlisting}

\section{Thiết kế kiến trúc triển khai và Hệ thống Route / API}

Hệ thống hỗ trợ chạy trực tiếp bằng Uvicorn hoặc đóng gói qua Docker Compose (\texttt{Dockerfile}, \texttt{docker-compose.yml}) ánh xạ cổng \texttt{8010:8000} kèm kiểm tra sức khỏe tại \texttt{/healthz} và \texttt{/readyz} (Hình~\ref{fig:g-deployment}). Bảng~\ref{tab:g-api-catalog} liệt kê các nhóm API chính trong \texttt{app/routes/api/}.

\begin{figure}[H]
\centering
\vspace{2pt}
\resizebox{0.98\textwidth}{!}{%
\begin{tikzpicture}[
    font=\sffamily\footnotesize,
    >=Latex,
    nodebox/.style={rectangle, draw=diagblueline, line width=0.9pt, fill=diagbluefill, rounded corners=5pt, text width=4.1cm, align=center, inner sep=6pt},
    arr/.style={-Latex, line width=0.9pt, diagblueline}
]

\node[nodebox] (client) {
\textbf{Máy trạm tại Xưởng}\\
Trình duyệt Chrome / Edge\\
React 19 SPA (\texttt{:3000} / \texttt{:5173})\\
Kết nối API \texttt{/api/*}
};

\node[nodebox, fill=diaggreenfill, draw=diaggreenline, right=1.8cm of client] (docker) {
\textbf{Máy chủ / Docker Compose}\\
Container \texttt{garageai-web}\\
Python 3.11 + Uvicorn (\texttt{8010:8000})\\
Healthcheck: \texttt{/healthz} \& \texttt{/readyz}
};

\node[nodebox, fill=diagslatefill, draw=diagslateline, right=1.8cm of docker, yshift=1.35cm] (storage) {
\textbf{Lưu trữ Bền vững (Disk)}\\
\texttt{database/garage.db} (WAL)\\
\texttt{backups/*.db} (\texttt{backup.ps1})\\
\texttt{config/prompts/*.txt}
};

\node[nodebox, fill=diagamberfill, draw=diagamberline, below=0.55cm of storage] (cloud) {
\textbf{Cổng Đám mây LLM API}\\
Groq (\texttt{openai/gpt-oss-120b})\\
Gemini / OpenRouter / OpenAI\\
HTTPS (TLS 1.3, Timeout 15s)
};

\draw[arr, <->] (client.east) -- node[above, font=\scriptsize]{HTTP + CSRF} (docker.west);
\draw[arr, <->, diagslateline] ([yshift=0.45cm]docker.east) -- (storage.west);
\draw[arr, <->, dashed, diagamberline] ([yshift=-0.45cm]docker.east) -- (cloud.west);

\end{tikzpicture}%
}

\vspace{2pt}
\caption{Sơ đồ triển khai hệ thống trên máy trạm tại xưởng hoặc Docker Compose}
\label{fig:g-deployment}
\end{figure}

\begin{table}[H]
\centering
\small
\caption{Danh mục các mô-đun JSON API chính trong \texttt{app/routes/api/}}
\label{tab:g-api-catalog}
\begin{tabularx}{\textwidth}{|L{4.0cm}|C{1.8cm}|L{3.1cm}|Y|}
\hline
\thead{Mô-đun \& Đường dẫn} & \thead{Phương thức} & \thead{Quyền kiểm tra} & \thead{Chức năng nghiệp vụ chính}\\
\hline
\texttt{auth\_api.py}\newline(\texttt{/api/auth/*}) & POST, GET & Công khai / Đã đăng nhập & Đăng nhập Argon2id, cấp cookie \texttt{garage\_session} + \texttt{garage\_csrf}, đổi mật khẩu (tăng \texttt{auth\_version}), đăng xuất. \\
\hline
\texttt{parties\_api.py}\newline(\texttt{/api/customers}, \texttt{/vehicles}, \texttt{/appointments}) & GET, POST, PUT, DELETE & \texttt{customer.*}, \texttt{vehicle.*}, \texttt{appointment.*} & CRUD khách hàng, phương tiện, lịch hẹn; tìm kiếm, lọc và phân trang. \\
\hline
\texttt{orders\_api.py}\newline(\texttt{/api/orders/*}) & GET, POST, DELETE & \texttt{order.*} + \texttt{duoc\_gan\_cong} & Tạo phiếu sửa chữa, lưu biên bản tiếp nhận (\texttt{intake\_records}), phân công KTV, thêm/gỡ hạng mục, đổi trạng thái, nghiệm thu. \\
\hline
\texttt{quotations\_api.py}\newline(\texttt{/api/quotations/*}) & GET, POST, DELETE & \texttt{quote.manage}, \texttt{quote.approve} & Tạo báo giá, thêm/gỡ dòng báo giá, gửi/duyệt báo giá và áp dụng \textit{idempotent} vào phiếu sửa chữa. \\
\hline
\texttt{catalog\_api.py} \& \texttt{stock\_api.py} & GET, POST, PUT & \texttt{catalog.*}, \texttt{stock.manage} & Quản lý danh mục dịch vụ, phụ tùng, nhập/điều chỉnh kho qua sổ cái \texttt{stock\_transactions} và đối soát tồn kho. \\
\hline
\texttt{billing\_api.py}\newline(\texttt{/api/invoices}, \texttt{/payments}) & GET, POST, DELETE & \texttt{invoice.manage}, \texttt{payment.record} & Lập hóa đơn cho phiếu đã \texttt{COMPLETED}, thu tiền từng phần/toàn bộ (\texttt{paid <= total}) và hủy khoản thu nhầm. \\
\hline
\texttt{warranties\_api.py}\newline(\texttt{/api/warranties/*}) & GET, POST & \texttt{warranty.manage} & Tạo bảo hành theo hạng mục của phiếu đã hoàn thành, yêu cầu bảo hành (\texttt{CLAIMED}) hoặc hủy bảo hành. \\
\hline
\texttt{ai\_api.py}\newline(\texttt{/api/ai/*}) & GET, POST & \texttt{ai.history}, \texttt{ai.explain}, \texttt{ai.quote}, \texttt{ai.review} & Chạy 3 chức năng AI01--AI03 và duyệt/từ chối kết quả AI trong bảng \texttt{ai\_results}. \\
\hline
\texttt{dashboard\_api.py}, \texttt{reports\_api.py}, \texttt{admin\_api.py} & GET, POST, PUT & \texttt{dashboard.view}, \texttt{report.view}, \texttt{user.*} & Tổng quan garage, 5 báo cáo thống kê và quản trị tài khoản nhân viên. \\
\hline
\end{tabularx}
\end{table}

\section{Kết luận chương}

Chương 4 đã trình bày đầy đủ thiết kế kiến trúc 4 tầng, lược đồ CSDL 21 bảng chuẩn 3NF với 37 chỉ mục, quy trình xử lý 7 bước của AI Gateway và 3 vòng thử nghiệm cải tiến Prompt. Chương 5 tiếp theo sẽ trình bày chi tiết việc hiện thực hóa các thiết kế này trong mã nguồn.
